\section{Proof of Theorem~\ref{thm:weighted-hoeffding-null}}
\label{app:hoeffding-null-proof}

We prove the theorem in four steps. First, we compute the second-order
expansion of the functional $D/30$ at the independence copula and identify
its quadratic kernel $K$. Second, we decompose the distinct-index statistic
into a weighted quadratic form in $K$ and a remainder that is negligible
under \eqref{eq:weight-diffuseness}. Third, we identify the limit law of the
quadratic form through the spectral decomposition of $K$. Finally, we derive
the variance expansion.

\begin{proof}
Apply the probability integral transform and write $Z_i=(U_i,V_i)$, where
all coordinates are independent uniform variables. Let $T(P)=D(P)/30$ and
let $P_0$ be the uniform measure on $[0,1]^2$.

\medskip\noindent
\emph{Step 1: second-order expansion of the functional.}
For $z=(u,v)$, contamination in the direction $\delta_z-P_0$ changes the
joint-minus-product discrepancy at $(x,y)$ at the rate
\begin{align*}
  \dot\Delta_z(x,y)
  &= \{\mathbf{1}(u\leq x)-x\}
     \{\mathbf{1}(v\leq y)-y\}.
\end{align*}
Because the discrepancy itself is zero at $P_0$, the first derivative of $T$
vanishes. The change in the integrating measure contributes only at third
order, while the mixed second derivative is
\begin{align*}
  K(z,z')
  &= \int_0^1\int_0^1
     \dot\Delta_z(x,y)\dot\Delta_{z'}(x,y)
     \,\mathrm{d}x\,\mathrm{d}y
   = k(u,u')k(v,v'),
\end{align*}
where direct integration gives
\begin{align*}
  k(a,b)
  &= \int_0^1
     \{\mathbf{1}(a\leq t)-t\}
     \{\mathbf{1}(b\leq t)-t\}\,\mathrm{d}t
   = \frac{a^2+b^2}{2}-\max(a,b)+\frac13.
\end{align*}
More precisely, let $\nu$ and $\gamma$ be zero-mass signed directions with
bounded densities relative to $P_0$, and set
$\dot\Delta_\nu(x,y)=\int \dot\Delta_z(x,y)\,\mathrm{d}\nu(z)$. Then
\begin{align*}
  \left.
  \frac{\partial^2}{\partial\varepsilon\,\partial\eta}
  T(P_0+\varepsilon\nu+\eta\gamma)
  \right|_{\varepsilon=\eta=0}
  &=2\int\!\!\int K(z,z')\,\mathrm{d}\nu(z)\,\mathrm{d}\gamma(z').
\end{align*}
The kernel $K$ is canonical in either argument. Therefore equality of these
bilinear forms over bounded zero-mean densities identifies the canonical
$L^2(P_0^2)$ quadratic-derivative kernel.

\medskip\noindent
\emph{Step 2: quadratic projection of the statistic.}
To connect this derivative with the distinct-index statistic, symmetrize the
order-$m$ kernels defining $\widehat\theta_{m,w}$ and denote their canonical
order-$r$ projections relative to $P_0$ by $h_{m,r}$. Differentiating
$T=\theta_3-2\theta_4+\theta_5$ once gives the first identity below. The
kernel representing its mixed second derivative is
$2\{3h_{3,2}-12h_{4,2}+10h_{5,2}\}$, and comparison of the two bilinear forms
for all $\nu$ and $\gamma$ as above gives the second identity in
$L^2(P_0^2)$:
\begin{align*}
  3h_{3,1}-8h_{4,1}+5h_{5,1}
  &=0,
  &
  3h_{3,2}-12h_{4,2}+10h_{5,2}
  &=K.
\end{align*}
The first identity proves degeneracy. The second identifies the full
quadratic projection and its normalization.

Put $c_{m,n}=m!e_m(p_1,\ldots,p_n)$ and, for an index set $A$,
$c_{m-r,n}^{(-A)}=(m-r)!e_{m-r}\bigl((p_i)_{i\notin A}\bigr)$. The canonical
decomposition of \citet{hoeffding1948u}, with the coefficient of each index
set replaced by its normalized distinct-tuple mass, gives
\begin{align*}
  \widehat\theta_{m,w}-\theta_m
  &= \sum_{r=1}^m
     \sum_{\substack{A\subseteq\{1,\ldots,n\}\\|A|=r}}
     \frac{
       (m)_r\left(\prod_{i\in A}p_i\right)
       c_{m-r,n}^{(-A)}
     }{c_{m,n}}
     h_{m,r}(Z_A),
\end{align*}
where $(m)_r=m!/(m-r)!$. Removing collisions from a fixed number of draws
gives, uniformly over all $A$ with $|A|\leq5$,
\begin{align*}
  |c_{m,n}-1|
  &\leq C_m s_{2,n},
  &
  |c_{m-r,n}^{(-A)}-1|
  &\leq C_m\left(s_{2,n}+\sum_{i\in A}p_i\right).
\end{align*}
Let $\alpha_3=1$, $\alpha_4=-2$, and $\alpha_5=1$, and define
\begin{align*}
  R_{1,n}
  &=\sum_i p_i\sum_{m=3}^5\alpha_m m
    \left\{
      \frac{c_{m-1,n}^{(-\{i\})}}{c_{m,n}}-1
    \right\}h_{m,1}(Z_i),\\
  R_{2,n}
  &=\sum_{i<l}p_ip_l\sum_{m=3}^5\alpha_m(m)_2
    \left\{
      \frac{c_{m-2,n}^{(-\{i,l\})}}{c_{m,n}}
      -\frac{1}{1-s_{2,n}}
    \right\}h_{m,2}(Z_i,Z_l),\\
  R_{\geq3,n}
  &=\sum_{m=3}^5\alpha_m
    \sum_{r=3}^m
    \sum_{\substack{A\subseteq\{1,\ldots,n\}\\|A|=r}}
    \frac{
      (m)_r\left(\prod_{i\in A}p_i\right)
      c_{m-r,n}^{(-A)}
    }{c_{m,n}}h_{m,r}(Z_A).
\end{align*}
Substitution of the two projection identities gives
\begin{align}
  \frac{\widehat D_w}{30}
  &= \frac{2}{1-s_{2,n}}
     \sum_{1\leq i<l\leq n}p_ip_lK(Z_i,Z_l)+R_n,
  \label{eq:hoeffding-quadratic-projection}
\end{align}
where $R_n=R_{1,n}+R_{2,n}+R_{\geq3,n}$. Canonical projections of different
orders or different index sets are orthogonal. The collision bounds, the
fixed $L^2(P_0^r)$ norms of the kernels, and
$e_r(p_1^2,\ldots,p_n^2)\leq s_{2,n}^r/r!$ therefore give
\begin{align*}
  \mathbb{E}_0(R_{1,n}^2\mid p)
  &\leq C s_{2,n}(\max_i p_i+s_{2,n})^2,\\
  \mathbb{E}_0(R_{2,n}^2\mid p)
  &\leq C s_{2,n}^2(\max_i p_i+s_{2,n})^2,\\
  \mathbb{E}_0(R_{\geq3,n}^2\mid p)
  &\leq C s_{2,n}^3.
\end{align*}
Consequently,
\begin{align*}
  \frac{\mathbb{E}_0(R_n^2\mid p)}{s_{2,n}^2}
  &\leq C\Biggl[
    \left\{
      \frac{\max_i p_i}{\sqrt{s_{2,n}}}+\sqrt{s_{2,n}}
    \right\}^2
    +(\max_i p_i+s_{2,n})^2+s_{2,n}
  \Biggr]
  \longrightarrow0,
\end{align*}
where the three terms come, respectively, from the residual first
projection, the quadratic coefficient error, and the projections of orders
three through five. The remainder is thus negligible at the scale $s_{2,n}$
of the quadratic term.

\medskip\noindent
\emph{Step 3: limit law of the quadratic projection.}
The one-dimensional kernel has the cosine expansion
\begin{align*}
  k(a,b)
  &= \sum_{j=1}^{\infty}
     \frac{2}{\pi^2j^2}\cos(\pi ja)\cos(\pi jb),
\end{align*}
so $K$ has orthonormal eigenfunctions and eigenvalues
\begin{align*}
  \phi_{jk}(u,v)
  &=2\cos(\pi ju)\cos(\pi kv),
  &
  \lambda_{jk}
  &=\frac{1}{\pi^4j^2k^2}.
\end{align*}
With $a_{i,n}=p_i/\sqrt{s_{2,n}}$, the leading term in
\eqref{eq:hoeffding-quadratic-projection}, multiplied by
$(1-s_{2,n})/s_{2,n}$, is
\begin{align*}
  \frac{2}{s_{2,n}}\sum_{i<l}p_ip_lK(Z_i,Z_l)
  &=
  \sum_{j,k\geq1}\lambda_{jk}
  \left[
    \left\{\sum_i a_{i,n}\phi_{jk}(Z_i)\right\}^2
    -\sum_i a_{i,n}^2\phi_{jk}(Z_i)^2
  \right].
\end{align*}
For any finite set of indices, the weighted Lindeberg central limit theorem
turns the first sums into independent standard normal variables. Since the
eigenfunctions are bounded and $\max_i a_{i,n}^2\to0$, the weighted law of
large numbers makes each second sum converge in probability to one.

It remains to pass from finite truncations to the full series. Let $K_M$
retain the eigenpairs with $1\leq j,k\leq M$. Canonicality gives the
uniform identity
\begin{align*}
  \mathbb{E}_0\!\left[
    \left|
      \frac{2}{s_{2,n}}\sum_{i<l}p_ip_l
      \{K-K_M\}(Z_i,Z_l)
    \right|^2
    \middle|p
  \right]
  &=
  2\left(1-\frac{s_{4,n}}{s_{2,n}^2}\right)
  \sum_{\substack{j>M\text{ or }k>M}}
  \lambda_{jk}^2\\
  &\leq
  2\sum_{\substack{j>M\text{ or }k>M}}
  \lambda_{jk}^2.
\end{align*}
The corresponding tail of the limiting centered Gaussian series has variance
twice the same eigenvalue-square sum. Since
$\sum_{j,k\geq1}\lambda_{jk}^2=1/8100<\infty$, both tail bounds vanish as
$M\to\infty$, which proves \eqref{eq:hoeffding-centered-limit}. The identity
$\sum_{j,k\geq1}\lambda_{jk}=1/36$ adds back the mean of the quadratic
series and proves \eqref{eq:hoeffding-positive-limit}.

\medskip\noindent
\emph{Step 4: variance expansion.}
The estimator is exactly unbiased because every normalized distinct tuple has
the same null expectation. Let $Q_n$ denote the quadratic term on the
right-hand side of \eqref{eq:hoeffding-quadratic-projection}. Canonicality
gives
\begin{align*}
  \operatorname{Var}_0(Q_n\mid p)
  &=\frac{s_{2,n}^2-s_{4,n}}
          {4050(1-s_{2,n})^2},
\end{align*}
so in particular $\operatorname{Var}_0(Q_n\mid p)=O(s_{2,n}^2)$. By
Cauchy--Schwarz,
\begin{align*}
  \left|
    \operatorname{Var}_0(Q_n+R_n\mid p)
    -\operatorname{Var}_0(Q_n\mid p)
  \right|
  &\leq
  \mathbb{E}_0(R_n^2\mid p)
  +2\sqrt{
    \operatorname{Var}_0(Q_n\mid p)
    \mathbb{E}_0(R_n^2\mid p)
  }
  =o(s_{2,n}^2).
\end{align*}
Multiplication by $30^2$ completes the variance calculation.
\end{proof}
